\cotefile{136}
\scanpage{297}
\[
\begin{aligned}
&\deleted{=\bigl(\underline{\mathrm{Hom}}^{*}((\tau_N\unsure{\mathcal{B}|B|S}))[-i],Q)\bigr)_0}\\
&\deleted{=\underline{\mathrm{Hom}}^{*}(\tau_{N-i}S)}\\
&\simeq\bigl(\underline{\mathrm{Hom}}^{*}(\tau_{N-i}(\mathcal{B}),\mathcal{A})\bigr)_i\\
&\simeq\underline{\mathrm{Hom}}^{i}(\tau_{N-i}(\mathcal{B}),\mathcal{A})
\end{aligned}
\]

Enfin si dans $C$ les $\underline{\lim}$ existent, alors $i_N$ admet un adjoint à gauche $q_N$.

En dualisant (appliquant ce qui précède à $C^o$, en utilisant l'isom
\diagram{\begin{tikzcd}[ampersand replacement=\&,column sep=small]
\mathcal{M}(C,S)^o \arrow[dd] \& \simeq \& \mathcal{M}(C^o,S)\arrow[dd] \\
\& \& \\
\mathcal{M}^{N}_{\deleted{\unsure{N|\leq N}}}(C,S)^o \& \simeq \& \mathcal{M}_{\leq N}(C^o,S)
\end{tikzcd}}
on voit que (si $C$ \unsure{ou} $C^o$ \unsure{est} une cat. de Groth. p.ex) le foncteur
\[
i^N:\mathcal{M}^N\longrightarrow\mathcal{M}
\]
\deleted{admet} \inserted{s'insère dans} une suite \deleted{\unsure{d'adjoints}} \inserted{de 4} foncteurs adjoints
\[
\varphi^{N}\ \tau^{N}\ i^{N}\ q^{N}
\]
qui se décrivent de façon duale \unsure{en|---}

% archive sheet number: 298
\scanpage{298}
\mspage{- 34 -}\hfill IX

On observera que $S'$ est noethérien régulier connexe de dimension 1, ce qui donne un sens à la conclusion énoncée. Pour prouver 3.6, choisissons une clôture séparable $\overline{K}$ de $K$, et écrivons l'extension $\overline{K}$ comme limite inductive de ses sous-extensions galoisiennes finies $K_i$. Soit $S_i$ le normalisé de $S$ dans $K_i$, et soit $U$ le plus grand ouvert de $S$ sur lequel $A_K$ ait bonne réduction. Comme $S$ est noethérien connexe de dimension 1, et $U\ne\varnothing$, on voit que $T=S-U$ est un ensemble fini de point fermés de $S$, et que pour tout $t\in T$, $\underline{O}_{S,t}$ est un anneau de valuation discrète. Il suffit de trouver, pour tout $t\in T$, un indice $i=i(t)$ tel que $A_{K_i}$ ait réduction semi-stable en les points de $S_i$ au-dessus de $t$ ; en effet, il suffira de prendre alors pour $K'$ le composé des extensions $K_{i(t)}$ pour $t\in T$. Cela nous ramène donc au cas où $S$ est un trait. D'autre part, comme de groupe $\mathrm{Gal}(K'/K)$ permute entre eux les points de $S_i$ au-dessus du point fermé $s$ de $S$, il suffit de trouver $K'$ de telle façon que $A_{K'}$ ait réduction semi-stable en \uline{un} point du normalisé $S'$ ; d'ailleurs il est inutile ici que $K'$ soit galoisien sur $K$, il suffit de trouver $K'$ extension finie étale de $K$ satisfaisant la condition précédente, car la clôture galoisienne de cette extension satisfera à la même condition. Utilisant le critère 3.5 (iv) on voit donc que 3.6 revient à l'assertion suivante :

\uline{Corollaire} 3.7. \uline{On suppose à nouveau} $S$ \uline{un trait}. \uline{Alors il existe un sous-groupe ouvert} $\handwritten{\pi}'$ \uline{de} $\handwritten{\pi}$ \uline{tel que l'action} $I'=I\handwritten{\cap\pi}'$ \uline{sur} $U=T_{\handwritten{\ell}}(A(\overline{K}))$ \uline{soit unipotente d'échelon} 2, \uline{ou encore} : \uline{il existe un sous-groupe ouvert} $I'$ \uline{de} $I$ \uline{dont l'action sur} $U$ \uline{soit unipotente d'échelon} 2.

\mspage{\handwritten{352}}
\mspage{352}
% archive sheet number: 299
\scanpage{299}
\[
\deleted{\unsure{\delta\leq d|\delta\geq d},\,\beta\geq0,}\quad \alpha,\beta\geq0,\ d\geq\beta-\alpha
\]
\diagram{\resizebox{\linewidth}{!}{\begin{tikzcd}[ampersand replacement=\&,column sep=small,row sep=large,cells={font=\scriptsize}]
\mathrm{Hom}^{d}(\tau_{\deleted{\unsure{\beta}}\alpha}M_{S_{\unsure{k|h}}},\tau_\beta N_{S_k})
 \arrow[r,"\sim","{\left(\substack{\alpha'\geq\alpha\\\beta'\geq\beta\\\deleted{\text{\unsure{restr.}}}\\\beta'-\alpha'\geq d}\right)}"']
 \& \mathrm{Hom}^{d}(\tau_{\deleted{\unsure{\beta}}\alpha'}M_{S_k},\tau_{\beta'}N_{S_k})
 \arrow[d,"\deleted{\unsure{1}}"',"{2}" description,"{\gamma\geq\deleted{\unsure{\alpha}}-d\;(\text{si }d\leq0)}"] \\
 \& \mathrm{Hom}^{d}(\tau_\gamma M,N)\quad
 \substack{\inserted{\searrow}\\(\text{sinon }\mathrm{Hom}^{d}(M,\tau_\gamma N)\\\text{si }0\leq\gamma\leq d,}
 \arrow[d,phantom,"{\rotatebox{90}{$\simeq$}}" description]\\
 \& \mathrm{Hom}^{d}(M_{S_k},\varphi_\beta\tau_\beta N_{S_k})\quad(\text{si }d\leq0)
 \arrow[d,phantom,"{\rotatebox{90}{$\simeq$}}" description]\\
 \& \mathrm{Hom}^{d}(S_k,\varphi_0(\mathrm{Hom}_{\deleted{\unsure{S_k}}k}(M,N)_{S_k}))
 \arrow[d,phantom,"{\rotatebox{90}{$\simeq$}}" description]\\
 \& \bigl(\varphi_0(\mathrm{Hom}_k(M,N)_{S_k})\bigr)_d
\end{tikzcd}}}

$M_{\deleted{\unsure{S|k}}}$ $k$-module
\[
 M_{S_k}=M\otimes_k k\{T\}
\]
\[
 (\varphi_0 M_{S_k})_{\unsure{-d|-i|d}}\simeq
 \underset{\substack{\rotatebox{90}{$\in$}\\u}}{\mathrm{Hom}^{-d}(\tau_d S_{\unsure{k}}, M_{S_k})}
\]
\[
 u(T^{(i)})=\xi_i T^{(i-d)}\qquad(i\geq d)\qquad \xi_i\in M
\]
\[
 u\bigl(\underbrace{T^{(i)}T^{(\ell)}}_{c_{i\ell}T^{(i+\ell)}}\bigr)
 =c_{i\ell}\xi_{i+\ell}T^{(i\unsure{-|+}\ell-d)}
\]
\[
 u(T^{(i)})T^{(\ell)}\mathbin{\unsure{:|=}}\xi_i T^{(i-d)}T^{(\ell)}
 =c_{i-d,\ell}\xi_i T^{(i\unsure{-|+}\ell-d)}
\]
\[
 \begin{aligned}
 i\geq d,\ \ell\geq0\quad &\Longrightarrow
 \boxed{c_{i,\ell}\xi_{i+\ell}=c_{i-d,\ell}\xi_i}\\[-2pt]
 &\hspace{9em}\unsure{\shortparallel}\\
 &\hspace{2em}c_{d,\ell}\xi_{d+\ell}=\xi_d
 \end{aligned}
\]
% archive sheet number: 300
\scanpage{300}
\mspage{- 160 -}\hfill IX

d'entiers $[1,r]$), $D_\lambda$ étant l'image par $\underline{f}:\underline{X}\handwritten{\longrightarrow}\underline{S}$ de la composante $D'_\lambda$ passant par le point singulier $x_\lambda$ de $X_{\unsure{o|0}}$ du lieu de non lissité $D'=\handwritten{\bigcup}_{\lambda\in L}D'_\lambda$ du morphisme $\underline{f}$. Par construction même, $f:X\handwritten{\longrightarrow} S$ est isomorphe à l'image inverse de $\underline{f}:\underline{X}\handwritten{\longrightarrow}\underline{S}$ par un $W$-morphisme convenable $g:S\handwritten{\longrightarrow}\underline{S}$ :

\diagram{\[
\begin{tikzcd}[ampersand replacement=\&]
 \underline{X} \arrow[d,"\underline{f}"'] \& X \arrow[l] \arrow[d,"f"]\\
 \underline{S} \& S \arrow[l,"g"'] \& .
\end{tikzcd}\tag*{(12.8.2)}
\]}

Posons
\[
 P=\underline{\mathrm{Pic}}^{o}_{\underline{X}/\underline{S}}\qquad ;\tag*{(12.8.3)}
\]
comme la fibre spéciale de $\underline{f}$ est géométriquement réduite de dimension 1, il s'ensuit que $\underline{F}$ est cohomologiquement plat en dimension 0 (EGA III 7.8.1) et que $P$ \deleted{\unsure{$-$}}est formellement lisse sur $S$ (ce dernier fait résultant de ce que $H^2(X_o,\underline{O}_{X_o})=0$, $X_o$ étant de dimension 1). Pour simplifier, nous allons admettre que $P$ \uline{est représentable par un schéma en groupes sur} $S$ (qui sera nécessairement un schéma en groupes lisse à fibres connexes) ; c'est là un cas particulier d'un résultat inédit de MUMFORD cité dans\handwritten{\unsure{\textsuperscript{1}}} [22](considérablement généralisé par RAYNAUD dans la note citée, dans le cas où la base est un trait). Comme aucune démonstration des résultats de MUMFORD ni de RAYNAUD n'a été publiée, le lecteur préfèrera peut-être utiliser les résultats généraux de M. ARTIN [1], qui donnent la représentabilité de $P$ (sous la seule hypothèse de partir d'un $\underline{f}$ propre, plat et cohomologiquement plat en dimension 0) par un ``espace algébrique''\unsure{.}

\mspage{478}
% archive sheet number: 301
\scanpage{301}
Notations $1/2$ simpliciales

Constructions universelles

% archive sheet number: 302
\scanpage{302}
\mspage{\textcircled{1}}
% slanted
\margin{1975 ou 1976}
\[
 \mathbb{\Delta}=\text{cat. des simplexes }\Delta^n\ (n\geq0)
 \qquad\substack{\Delta^n=[0,n]\subset\mathbb{N}\\\text{avec rel. d'ordre total}}
\]
\[
 \widehat{\mathbb{\Delta}}=\mathrm{Ss}
 =\underline{\mathrm{Hom}}(\mathbb{\Delta}^o,\mathrm{E}\mathrm{ns}))
 =\mathbb{E}\mathrm{ns}_*
\]

Plus généralement, si $A$ est une catégorie, on pose
\[
 \begin{aligned}
 A_*&=\underline{\mathrm{Hom}}(\mathbb{\Delta}^o,\mathbb{E}\mathrm{ns})\\
 A^*&=\underline{\mathrm{Hom}}(\mathbb{\Delta},\mathrm{Ens})
 \end{aligned}
\]
donc
\[
 (A^o)^*=(A_*)^o\qquad A^*\simeq((A^o)_*)^o
\]
où l'exposant $o$ désigne le \inserted{passage à la} cat. opposée.

Les objets de $A_*$ (resp. $A^*$) sont considérés comme des structures algébriques de type simple (sur une infinité d'objets de base) dans $A$, les ``objets $1/2$ simpliciaux'' resp ``$1/2$ cosimpliciaux''. \inserted{On} écrira $\mathbb{E}\mathrm{ns}_*$ quand on a en vue cet aspect, et $\widehat{\mathbb{\Delta}}$ \inserted{ou $\mathrm{Ss}$} quand on a plutôt le pt de vue ``objet d'un topos'' ou ``faisceaux sur un topos''\deleted{\unsure{,}}\textsuperscript{$\circ\!/$}% 
% margin note written sideways
\margin{\textsuperscript{$\circ\!/$} La notation $\widehat{\mathbb{\Delta}}$ peut être utilisée de préf. quand on veut garder à l'esprit la nature algébrique très particulière de ce topos, la notation $\mathrm{Ss}$ quand \deleted{\unsure{on}} veut l'oublier au profit de sa signification topologique.}
Tous les développements qui suivent ont pour objet d'étudier ces objets (l'équivalent combinatoire des espaces topologiques) et leurs ``types d'homotopie'', via des invariants qu'on peut leur associer, qui en dépendent de façon covariante ou contravariante. Nous avons ici en vue une étude systématique, dans l'esprit de l'algèbre universelle, de ces invariants, (qui sont donc des foncteurs sur $\widehat{\mathbb{\Delta}}\unsure{\simeq|=}\mathbb{E}\mathrm{ns}_*$) \deleted{et des} \inserted{de leurs structures et des} opérations qui peuvent être définies sur elles, permettant d'en déduire certains complexes \unsure{:} à partir d'autres plus simples

% archive sheet number: 303
\scanpage{303}
\mspage{\textcircled{2}}
Un objet de $A_*$ sera généralement noté par un symbole de la forme $K_*$, où $K_*$ désigne la famille des
\[
 K_*(\Delta^n)=K_{[n]}\qquad
 K_*\overset{\inserted{\text{par abus de notation}}}{=}(K_{[n]})_{n\geq0}
 \quad\inserted{\text{ou mieux }\Delta^n\longmapsto K_{[n]}}
\]
avec les opérations $1/2$ simpliciales entre elles. On fera attention qu'on écrit $K_{[n]}$ et non $K_n$, pour des raisons qui vont apparaître (\unsure{impérieuses|limpidement} lorsque $A$ est additive\ldots). De $=$ un objet de $A^*$ sera noté
\[
 K^*=(n\longmapsto K^{[n]})\overset{\text{abrév.}}{=}(K^{[n]})_n
\]

Nous aurons à travailler avec la situation où on a deux catégories $A$ et $B$, et une équivalence (notée \deleted{$\unsure{*}$} $M\longmapsto M^\vee$)
\[
 v:A^o\xrightarrow{\approx}B\qquad\text{d'où}\quad B^o\overset{\approx}{\unsure{\longrightarrow|-}}A
\]
[le plus souvent $A=B$ et $(\deleted{\unsure{M|\check{M}}}M^\vee)^\vee\underset{\substack{\text{isom}\\\text{fonct.}}}{\simeq}M$, avec compatibilités habituelles d'une autodualité\ldots]
on notera alors \inserted{souvent} par la $=$ lettre un objet de $A_*$ (ou $A^*$) et l'objet de $A^*$ (resp $A_*$) qui lui correspond par application de $v$, mais en indiquant la variance par la position du signe $*$
\[
 \begin{aligned}
 K^*&=(K_*)^\vee & K^{[n]}&=K_{[n]}^\vee\\
 K_*&=(K^*)^\vee & K_{[n]}&=K^{[n]\,\vee}
 \end{aligned}
\]

On s'intéressera surtout au cas où ($k$ étant un anneau fixé) \deleted{\unsure{$A=\text{la}$}}
\[
 \begin{aligned}
 \deleted{\text{cat. }A}\qquad A&=\deleted{\unsure{\mathrm{Ab}_k}}\ k\text{-modules à g. proj. de t.f.}\\
 B&=k\text{-modules à droite proj. de t.f.}
 \end{aligned}
\]
(si $k$ est commutatif, on a $A=B$ et on a une autodualité \unsure{sans|$v$ avec|dans|sur} $A$).

On a le foncteur canonique pl. fid
\[
 \mathbb{\Delta}\lhook\joinrel\longrightarrow\widehat{\mathbb{\Delta}}=\mathrm{Ss}\tag*{(*)}
\]
l'image par ce foncteur de $\Delta^n$ est noté $\Delta^n_*$ \inserted{ou mieux $\Delta^{[n]}_*$ (puisqu'il est covariant en $n$)}. On a donc
\[
 \Delta^n_*\ (\text{ou }\Delta^{[n]}_*)=
 \bigl(m\longmapsto\Delta^n_{[m]}=\mathrm{Hom}_{\mathrm{Ens}}(\Delta^m,\Delta^n)\bigr)
\]
et pour $\deleted{\unsure{F}}X_*\in\mathrm{Ob}\,\widehat{\mathbb{\Delta}}$,
\[
 \mathrm{Hom}(\Delta^n_*,\deleted{\unsure{F}}X_*)\simeq X_*(\Delta^n)\simeq X_{[n]}
\]
en particulier
\[
 \mathrm{Hom}(\Delta^n_*,\Delta^m_*)\simeq\Delta^m_*(\Delta^n_*)\simeq
 \underset{\substack{\Vert\\(\mathrm{Hom}_{\mathrm{ord}}(\Delta^n,\Delta^m))}}{\mathrm{Hom}_{\mathbb{\Delta}}(\Delta^n,\Delta^m)}
\]
Pour $n$ variable, $\Delta^n_*\in\widehat{\mathbb{\Delta}}$ dépend de façon covariante de $n$, \deleted{\unsure{i.e. o.|i.e.}} donc le foncteur $n\longmapsto\Delta^n_*$ (qui n'est autre que $(*)$) est un objet canonique de $\widehat{\mathbb{\Delta}}^*$, qu'on appelle \uline{objet cosimplicial universel} et note $\Delta^*$ (ou mieux $\Delta^*_* $)
\[
 \begin{aligned}
 \Delta^*_*\in\widehat{\mathbb{\Delta}}^* &=(\mathrm{Ens}_*)^*
 \quad\text{d'où }\forall n\geq0\\
 \Delta^{[n]}_*\in\widehat{\mathbb{\Delta}} \qquad
 \Delta^{[n]}_*&=\Delta^n_* =m\longmapsto
 \underset{\substack{\Vert\\\mathrm{Hom}_{\mathbb{\Delta}}(\Delta^m,\Delta^{\unsure{m|n}})}}{\Delta^n_{[m]}}
 \end{aligned}
\]

\deleted{\unsure{Ainsi}} [Une parenthèse de notations : Considérons une catégorie de la forme $(A_*)^*$, un objet de la dite est donc noté $K^*=\unsure{(|[}K^{[n]})_n$, où les $K^{[n]}$ sont des $\in A_*$, donc s'écrivent $K^{[n]}_*$,
\[
 K^{[n]}_*=(K^{[n]}_{[m]})_{m\geq0}
\]
Ainsi, on en vient à noter
\[
 K^*_*\quad\text{les objets de }(A_*)^*\simeq
 \underline{\mathrm{Hom}}(\mathbb{\Delta}\times\mathbb{\Delta}^o,A\unsure{)}
\]
\deleted{\unsure{on}}
% archive sheet number: 304
\scanpage{304}
\mspage{\textcircled{3}}
\[
 K^{[m]}{}_{[n]}=K(\Delta^m,\Delta^n)\in A
\]
De $=$, on note
\[
 \begin{aligned}
 K_*{}^*&\quad\text{les objets de }(A^*)_*\simeq\underline{\mathrm{Hom}}(\mathbb{\Delta}^o\times\mathbb{\Delta},A)\\
 K^{**}&\quad\text{les objets de }(A^*)^*\simeq\underline{\mathrm{Hom}}(\mathbb{\Delta}^{\unsure{o}}\times\mathbb{\Delta},A)\\
 K_{**}&\quad\text{les objets de }A_{**}=\underline{\mathrm{Hom}}(\mathbb{\Delta}^o\times\mathbb{\Delta}^o,A)
 \end{aligned}
\]

On fera attention que si on compose avec \inserted{l'un} des foncteurs de symétrie \deleted{$\unsure{(m,n)\mapsto(n,m)}$} \inserted{$(\Delta^m,\Delta^n)\longmapsto(\Delta^n,\Delta^m)$} dans $\mathbb{\Delta}\times\mathbb{\Delta}$ ou dans $\mathbb{\Delta}^o\times\mathbb{\Delta}^o$, ou entre $\mathbb{\Delta}\times\mathbb{\Delta}^o$ et $\mathbb{\Delta}^o\times\mathbb{\Delta}$), on n'utilise pas la $=$ lettre pour désigner un \deleted{complexe} \inserted{objet} (tel que $K^*{}_*$ ou $K^{**}$ disons) ou son transformé (qui devrait alors s'écrire $K_*{}^*$, ou encore une fois $K^{**}$), car cela conflicterait avec la convention plus haut de la ``montée et descente des indices'' $K_*{}^*\simeq K^*{}_*{}^\vee$ dans le premier cas, et \deleted{\unsure{serait}} amènerait à écrire $K^{[m][n]}=K^{[n][m]}$ dans le \unsure{deuxième} ! On utilisera au besoin un symbole ``de symétrie'' tel que $\sigma$, pour dénoter le symétrique de $K$ ($=K^*{}_*,K^{**}$) par $ {}^{\sigma}K$ ($=({}^{\sigma}K)_*{}^*,({}^{\sigma}K)^{**}$) \deleted{\unsure{de|dans}}\ldots\ D'ailleurs, \deleted{Ces considérations et notations} si on a un objet $K^*{}_*\inserted{{}\in\mathrm{Ob}\,(A_{\unsure{*}}^{*})^{*}_{\deleted{\unsure{*}}}}$ disons, alors pour \unsure{tt} $n\geq0$, on peut considérer
\[
 \deleted{\unsure{K\longmapsto}}\quad K^*{}_{[n]}\overset{\text{déf}}{=}(m\longmapsto K^{[m]}{}_{[n]})
\]
comme un objet de $A^*$ ; le foncteur $[n]\longmapsto K^*{}_{[n]}$ de $A^*{}_*$ n'est autre que $ {}^{\sigma}K$.

On étend ces considérations et notations au cas de symboles $K$ à plus de deux indices, pour désigner des objets de catégories du
% archive sheet number: 305
\scanpage{305}
\mspage{\textcircled{4}}
la forme $A^*{}_*{}^{**}=(((A^*)_*)^*)^*$ \deleted{\unsure{etc, et de}} p.ex. -- dont les objets sont notés \deleted{par} par des lettres à étoile $K^{*}\unsure{{}^{*}_{*}}{}_*{}^{*}$ \ldots\ ]

Le rôle de $\Delta^*_{\deleted{\unsure{*}}}\in\mathrm{Ob}\,\widehat{\mathbb{\Delta}}^*\inserted{{}=\mathrm{Ob}\,\mathrm{Ss}^*}$ comme ``objet cosimplicial universel'' est explicité par la

\uline{Proposition} 1. Soit $A$ une catégorie où les $\underline{\lim}$ existent (resp $B$ une catégorie où les $\varprojlim$ existent). Alors
\[
 \begin{aligned}
 \text{a)}\quad A^*&\underset{\leftarrow}{\approx}\underline{\mathrm{Hom}}_{\substack{\text{comm}\\\text{à }\varinjlim}}(\deleted{\unsure{\widehat{\mathbb{\Delta}}}}\mathrm{Ss},A)\\
 \text{par}\quad F&\longmapsto F(\Delta^*)\\
 \text{b)}\quad B_*&\underset{\leftarrow}{\approx}\underline{\mathrm{Hom}}_{\substack{\text{comm}\\\text{à }\varprojlim}}(\deleted{\unsure{\widehat{\mathbb{\Delta}}}}\mathrm{Ss}^{o},\deleted{\unsure{A}}B)\\
 \text{par}\quad F&\longmapsto F((\Delta^*)^o)
 \end{aligned}
\]
(où $(\Delta^*)^o$ désigne l'objet simplicial de $\widehat{\mathbb{\Delta}}^o$ défini par l'objet cosimplicial $\Delta^*$ de $\widehat{\mathbb{\Delta}}$).

Bien sûr, b) n'est autre que a) appliqué à $A=B^o$, et a) signifie que le foncteur ``restriction à $\mathbb{\Delta}$''
\[
 \underline{\mathrm{Hom}}(\widehat{\mathbb{\Delta}},A)\longrightarrow\underline{\mathrm{Hom}}(\mathbb{\Delta},A)
\]
induit une équivalence sur la \unsure{ss-cat|s-cat.} \deleted{\unsure{des fonct|$\underline{\mathrm{Hom}}$}} \inserted{comm.\ à $\underline{\lim}$} du premier \unsure{mb.|membre} C'est vrai pour \uline{\unsure{tte|toute|les}} \deleted{(petite)} \inserted{\unsure{mt|toute}} catégorie $C$ -- pas seulement pour $\mathbb{\Delta}$, et bien connu. Le foncteur en sens inverse associe, à $K_*=\varphi:\mathbb{\Delta}\unsure{\longrightarrow|-} A$, le foncteur
\[
 X_*\longmapsto
 \underline{\lim}_{\Delta^n_*\in\mathbb{\Delta}/X_*}\varphi(\Delta^n_*)
 =\underline{\lim}_{\Delta^n_*\in\mathbb{\Delta}/X_*}K_{[n]}
\]
Dans le cas b) si on écrit
\[
 \underline{\mathrm{Hom}}(\widehat{\mathbb{\Delta}}^o,\deleted{\unsure{A}}B)\longrightarrow
 \underline{\mathrm{Hom}}(\mathbb{\Delta}^o,B)
\]
le foncteur en sens inverse associe à
\[
 K^*=\psi:\mathbb{\Delta}^o\unsure{\longrightarrow|-}B
\]
le foncteur contravariant
\[
 X^*\longmapsto\varprojlim_{\Delta^n_*\in(\mathbb{\Delta}/X_*)^o}\varphi(\Delta^n)
 =\varprojlim_{\Delta^n\in(\mathbb{\Delta}/X_*)^o}K^{[n]}
\]

\uline{Yoga} : nous nous intéresserons surtout \deleted{aux} (en première étape) aux invariants covariants attachés à $X_*$, à valeurs dans une cat. $A$, qui \deleted{\unsure{pour $X_*$ être|commutent aux $X_*$}} \inserted{commutent} aux $\underline{\lim}$ p.r. à $X_*$, -- ils sont \deleted{\unsure{donc}} exprimés (si $A$ est ``grande'' pour contenir les $\underline{\lim}$ \unsure{qqc|qqles|voulues}) par des éléments de $A^*$ i.e. des objets \uline{cosimpliciaux} de $A$. Pour les invariants contravariants en $X_*$, à val. dans une $B$, \deleted{\unsure{alors|dans}} nous attacherons en premier lieu à ceux qui \deleted{\unsure{con commutent}} aux $\varprojlim$ en $X_*$ ; pour $B$ stable par $\varprojlim$, ils correspondent \deleted{donc} aux objets de $B_*$, i.e. les objets $1/2$ simpliciaux de $B$. Si $B=(\mathrm{Ens})$, ils correspondent donc à des ens. $1/2$ simpliciaux, i.e. à des objets $K_*$ de $\deleted{\unsure{\widehat{\mathbb{\Delta}}}}\mathrm{Ss}$ lui-$=$.
% archive sheet number: 306
\scanpage{306}
\mspage{5 \deleted{\unsure{5}}}
Ce n'est autre que le foncteur représenté par $K_*$, et on le notera aussi
\[
 \deleted{\unsure{K_*(X_*)\longmapsto|K_*(X_*)=}}\qquad
 X_*\longmapsto K_*(X_*)=\mathrm{Hom}(X_*,K_*)
\]
Si $B=(\mathrm{Ab})$, $K_*$ sera de $=$ un objet de $\mathrm{Ss}_{\mathrm{ab}}$ \deleted{\unsure{$\mathrm{Ab}_*$}}, et en effet le foncteur qu'il représente est automatiquement muni de \deleted{\unsure{telle str.}} \inserted{\unsure{'une structure|type structure|une structure}} \unsure{abélien|abélienne}, $\overline{M}$ remarque pour tte structure algébrique (sur un ou plusieurs objets de base) qui ``puisse s'exprimer en termes de $\varprojlim$ exclusivement'' (grps, \deleted{\unsure{unités|modules}} anneaux, modules sur \deleted{\unsure{des algèbres, etc|tels, etc}} etc).

\uline{Rem} Les \deleted{foncteurs} invariants ainsi obtenus à coups de foncteurs représentables -- plus gén, d'objets simpliciaux ou cosimpliciaux de catégories $B$ ou $A$ -- sont de nature trop ``grosse'' et ``fruste'' pour être intéressants directement -- \deleted{\unsure{en}} en particulier, ce ne sont pas des invariants du type d'homotopie de $X_*$. On devra en extraire des invariants plus subtils qui soient des invariants du type d'homotopie de $X_*$, -- ceux-ci ne seront pas de nature si simple -- i.e. ``repr.'' par des objets simpliciaux ou cosimpliciaux. Notre propos sera de mettre sur les invariants ``gros'' qu'on vient de décrire, un maximum de structures supplémentaires \unsure{raisonnables}, et dans les catégories de structures de ce type (complexes de chaines ou de cochaines, algèbres cosimpliciales etc) d'inverser les flèches qu'on obtient à partir d'équiv. d'homotopie des $\mathrm{Ss}$, et de passer à des ``catégories de fractions'' en inversant ces flèches. L'invariant ainsi obtenu \deleted{\unsure{donne certes}} \inserted{(plus ``grossier''} bien sûr, mais plus subtil dans sa définition, et moins ``redondant'') exprimera alors de façon plus ou moins complète le type d'homotopie, ou en saisira de façon adéquate tels ou tels aspects.

Dans le cas où $A$ est une catégorie additive karoubienne (tout projecteur sur un objet a un noyau, ou ce qui revient au $=$, tt proj. a un conoyau) le \unsure{topoïs|topos|topo|théorie} de Dold-Puppe donne des équivalences canoniques
\diagram{\begin{tikzcd}[ampersand replacement=\&]
 A_* \arrow[r,shift left=1.1ex,"{\substack{\mathrm{ND}_{.}\\\approx}}"]
 \& A_{.}\arrow[l,shift left=1.1ex,"{\approx}"',"{\mathrm{DP}_*}"] \& {,}
\end{tikzcd}}
\margin{$\mathrm{ND}=$ ``non dégénéré''\\
$\mathrm{DP}=$ ``Dold-Puppe''}
où $A_.$ désigne la catégorie des \uline{complexes de chaines} dans $A$. On désigne par une $=$ lettre \inserted{telle $K$} deux objets qui se correspondent, mais en mettant en indice soit $*$ soit $.$, suivant qu'il
% archive sheet number: 307
\scanpage{307}
\mspage{(6)}
s'agit de désigner l'avatar ``objet simplicial'' ou ``complexe de chaines'' :
\[
 \begin{aligned}
 K_.&=\mathrm{ND}_ .(K_*), &K_*&=\mathrm{DP}_*(K_.)\\
 K_*&=(n\unsure{\mapsto|-}K_{[n]}), &K_.&=(K_n,d_n)_{n\geq0}
 \end{aligned}
\]
On est ainsi amené à noter par une notation $K_.$ \unsure{tt|tel} complexe de chaines, par $K_n$ ses composantes (ce qui exclut l'utilisation de la notation $K_n$ pour la $n$.\unsure{ème|ième} composante \inserted{de} $K_*$, et nous oblige à adopter une notation telle que $K_{[n]}$ !).

De $=$, le \unsure{topoïs|topos|topo|théorie} de DP appliqué à $A^o$ nous donne des équivalences
\diagram{\begin{tikzcd}[ampersand replacement=\&]
 A^* \arrow[r,shift left=1.2ex,"{\substack{\mathrm{ND}^{.}\\\approx}}"]
 \& A^{.}\arrow[l,shift left=1.2ex,"{\approx}"',"{\mathrm{DP}^{*}}"]
\end{tikzcd}}
et des notations
\[
 \begin{aligned}
 K^{.}&=\mathrm{ND}^{.}(K^*), & K^*&=\mathrm{DP}^*(K^{.})\\
 K_{\deleted{\unsure{*}}}^{*}&=(n\mapsto K^{[n]}), &K^{.}&=(K^n,d^n)_{n\geq0}.
 \end{aligned}
\]

Les foncteurs $\mathrm{ND}^{.}$ et $\mathrm{DP}^*$ commutent aux foncteurs additifs. En particulier, dans le cas d'une équivalence
\[
 A\underset{v}{\overset{\approx}{\longleftrightarrow}}B^o
\]
déjà \unsure{envisagé|envisagée}, étendant aux complexes de chaines ou de cochaines la \unsure{convent.|conv.|convention} de la montée ou de la descente des indices, on \unsure{trouve} que les foncteurs ND et DP y commutent. Donc dans ce cas, la $=$ lettre $K$ désigne
% archive sheet number: 308
\scanpage{308}
\mspage{7}
quatre avatars possibles
\[
 K^*\qquad K_*\qquad K^{.}\qquad K_.
\]
liés par les relations ci dessus
\[
 \left\{\begin{aligned}
 K^*&=\mathrm{DP}^*(K^{.}) & K^{.}&=\mathrm{ND}^{.}(K^*)\\
 K_*&=\mathrm{DP}_*(K_.) & K_.&=\mathrm{ND}_ .(K_*)
 \end{aligned}\right.
\]
plus les relations
\[
 \left\{\begin{aligned}
 K^*&=(K_*)^\vee, & K_*&=(K^*)^\vee\\
 K^{.}&=(K_.)^\vee, & K_.&=(K^{.})^\vee .
 \end{aligned}\right.
\]
\deleted{\unsure{(et notre remarque tient}}

Ces \deleted{$\unsure{K^*=|K}$} \unsure{notationss|notations} et remarques (i.e. équivalences de catégories) s'étendent aux objets \deleted{\unsure{dépendant}} \inserted{dépendant de} plusieurs indices, p.ex \deleted{$K^*$} une même lettre $K$ (désignant \inserted{\unsure{initialemt|p.ex. initialement}} un objet $K^*{}_*$ de $(A_*)^*$ \inserted{disons)}, désigne également, \deleted{\unsure{au total juste}} quatre avatars distincts
\[
 \left\{\begin{aligned}
 K^*{}_*&\in(A_*)^*, & (K^{.}{}_*)&\in(A_*^{\deleted{\unsure{*}}})^{.}\\
 K^{.}{}_.&\in(A_.)^{.}, &(K^*{}_.)&\in(A_.)^*
 \end{aligned}\right.
\]
avec des notations
\[
 \begin{aligned}
 K^{[n]}{}_{[m]},&\qquad K^n{}_{[m]}\\
 K^n{}_{m},&\qquad K^{[n]}{}_{m}
 \end{aligned}
\]
pour les composantes. De $=$, on aura, si on part p.ex. de $K^{**}\in(A^*)^*$
\[
 \left\{\begin{aligned}
 K^{**}&\in(A^*)^*, &K^{\unsure{.}*}&\in(A^{.})^*\\
 K^{..}&\in(A^{.})^{.}, &K^{*.}&\in(A^*)^{.}
 \end{aligned}\right.
\]
\[
 \left\{\begin{aligned}
 K^{[n][m]},&\qquad K^{n[m]}\\
 K^{nm},&\qquad K^{[n]m}
 \end{aligned}\right.
\]
On notera qu'un objet de $(A_*)^{.}$ est un complexe de cochaines \inserted{$n\longmapsto(K_*{}^n,d^n)$} dans la catégorie $A_*$ des objets $1/2$ simpliciaux de $A$,
% archive sheet number: 309
\scanpage{309}
\mspage{(8)}
ou au choix (moyennant \deleted{\unsure{une}} l'équivalence de catégories canonique
\[
 \sigma:(A_*)^{.}\underset{=}{\approx}(A^{.})_*\qquad)
\]
un objet $1/2$ simplicial dans la catégorie $A^{.}$ des complexes de cochaines, \unsure{noté|soit}
\[
 n\longmapsto K_{[n]}{}^{.}.
\]
On interprète de $=$ les objets des autres catégories mixtes \deleted{$\unsure{1/2\text{ simpliciales}}$}.
\[
 \left\{\begin{aligned}
 &A_{*.},\, A_{.*},\, A^*{}_{\,.}{}^{\deleted{\unsure{*}}},\,\deleted{\unsure{A^{.}{}_*}}\\
 &A^{*.}_{\deleted{\unsure{*}}},\, A^{.*}_{\deleted{\unsure{*}}},\, A^{.}{}_*\,\deleted{\unsure{A_*{}^{.}}}.
 \end{aligned}\right.\tag*{(*)}
\]
De façon analogue, les objets $K^{..}$ de $(A^{.})^{.}$ sont les bicomplexes de cochaines, les objets $K_{..}$ de $(A_.)_.=A_{..}$ sont les bicomplexes de chaines, enfin les objets $K^{.}{}_.$ de $(A_.)^{.}$ ou $K_.{}^{.}$ de $(A^{.})_.$ sont les bicomplexes mixtes \inserted{chaines-cochaines} (de chaines pour un indice, de cochaines pour un autre). Dans le cas de $K^{..}$ ou $K_{..}$ les complexes simples \unsure{associé|associés} est un complexe de cochaines resp de chaines, tandis que dans le cas mixte $K^{.}{}_.$ ou $K_.{}^{.}$, on trouve des complexes simples avec des objets en degrés positifs ou négatifs à priori (Il y aurait lieu de fixer une convention de degré, si pour le degré total l'op \unsure{diff|diff.} est de degré $+1$ ou $-1$, suivant
% archive sheet number: 310
\scanpage{310}
\mspage{(9)}
qu'on donne la préférence \unsure{aux}\deleted{\unsure{l'un|chaines}} cochaines ou aux chaines. (Je n'ai pas eu un besoin urgent d'une telle convention de notation\ldots)

Dans le cas d'une dualité
\[
 v:A\overset{\approx}{\rightleftarrows}B^o
\]
on étend les conventions de montée et descente des indices aux objets à plusieurs indices, ainsi
\[
 (K_*{}^{.})^\vee=K^*{}_.\ \deleted{\unsure{\}}}.
\]
Dans la suite d'objets mixtes $(*)$ de la page précédente, on a écrit l'un en dessous de l'autre des objets qui se correspondent par dualité.

\uline{Remarque} Il est essentiel pour notre système de conventions de définir (disons) des bicomplexes de cochaines (objets de $A^{..}$) avec deux op \unsure{diff.|de diff.} $d^{\deleted{\unsure{(1)}}\, .}_{(1)}$ et $d^{.}_{(2)}$ qui \uline{commutent} -- on n'utilisera donc pas les conventions de Cartan-Eilenberg. Il faudra donc un signe à la diff. totale.

Du point de vue de la prop. 1 et du yoga \unsure{correspondant} ``description d'invariants associés à des $X_*\in\mathrm{Ob}\,\mathrm{Ss}$'', on peut dire que le \unsure{topoïs|topos|topo|théorie} de Dold-Puppe
% archive sheet number: 311
\scanpage{311}
\mspage{10}
% archive sheet number: 312
fournit \deleted{un}

\uline{Prop. 2} Avec les notations \inserted{et hyp.} de prop. 1, avec $A$ et $B$ additives \deleted{karoubiennes}

\[
\begin{aligned}
\text{a)}\quad&\underline{\operatorname{Hom}}_{\substack{\text{comm}\\\text{à }\varinjlim}}(\mathrm{Ss},A)\xrightarrow{\sim} A^{.}\\
\text{b)}\quad&\underline{\operatorname{Hom}}_{\substack{\text{comm}\\\text{à }\varprojlim}}((\mathrm{Ss})^o,B)\xrightarrow{\sim} B_{.}
\end{aligned}
\]

Nous allons \deleted{\unsure{préciser}}\inserted{expliciter} cette correspondance du point de vue de l'algèbre universelle. Prenons d'abord le cas de b), avec $B=(\mathrm{Ab})$ pour simplifier. \deleted{\unsure{résultat|résulte}}\inserted{Donc il faut} préciser la relation entre
\[
F:(\mathrm{Ss})^o\xrightarrow[\substack{\text{commute}\\\text{aux }\varprojlim}]{}\mathrm{Ab} \qquad\text{et}\quad K_{.}=(K_n,d_n)\in \mathrm{Ab}_{.}
\]
qui se déterminent mutuellement. On désigne comme de juste par $K_*$ l'objet de $\mathrm{Ab}_*$ défini par $K_{.}$. On a donc
\[
F(X_*)\simeq\operatorname{Hom}_{\widehat{\mathbb{\Delta}}}(X_*,K_*)\simeq\operatorname{Hom}_{\widehat{\mathbb{\Delta}}_{ab}}(\mathbb{Z}(X_*),K_*)
\]
où $\deleted{\unsure{\{|(}}\mathbb{Z}(X_*)\deleted{\unsure{\}|)}}\deleted{\unsure{{}_*|*}}$ désigne le ``faisceau abélien libre engendré par $X_*$'' en l'occurrence
\[
\mathbb{Z}(X_*)=(n\longmapsto\mathbb{Z}(X_n)=\mathbb{Z}^{(X_n)})
\]
Ce groupe \inserted{ab} $1/2$ simplicial sera aussi noté $\mathbb{C}_*(X_*)$, et $\mathbb{C}_{.}(X_*)$ le complexe de chaines associé \inserted{par DP} : \deleted{On aura donc}
\[
\left\{\begin{aligned}
\mathbb{C}_*(X_*)&=\mathbb{Z}(X_*)=(n\longmapsto\mathbb{Z}(X_n)=\mathbb{Z}^{(X_n)}\\
\mathbb{C}_{.}(X_*)&=\mathrm{ND}_{.}(\mathbb{C}_*(X_*))
\end{aligned}\right.
\]
Donc par Dold-Puppe on aura
\[
F(X_*)\simeq\operatorname{Hom}_{\substack{\widehat{\mathbb{\Delta}}_{ab}\\\shortparallel\\\mathrm{Ab}_*}}(\mathbb{C}_*(X_*),K_*)\underset{\text{par D.P.}}{\simeq}\operatorname{Hom}_{\mathrm{Ab}_{.}}(\mathbb{C}_{.}(X_*),K_{.})
\]
Donc le foncteur $F_{K_{.}}=F\quad \mathrm{Ss}^o\longrightarrow \mathrm{Ab}$
\[
X_*\longmapsto F_{K_{.}}(X_*)=F(X_*)
\]
associé à $K_{.}\in \mathrm{Ab}_{.}$ est donné par
\[
\boxed{F_{\unsure{K_{.}|K_*}}(X_*)=\operatorname{Hom}_{.}(\mathbb{C}_{.}(X_*),K_{.})}
\]
On en tire deux conclusions :

1°) Les valeurs sur $X_*$ des foncteurs envisagés $F$ ($\mathrm{Ss}^o\longrightarrow \mathrm{Ab}$, commutant aux $\varprojlim$) ne dépendent que de $\mathbb{C}_*(X_*)\in \mathrm{Ab}_*$ \inserted{(``abélianisé'' de $X_*$)} --- ce qui revient au $=$, de $\mathbb{C}_{.}(X_*)\in \mathrm{Ab}_{.}$ ; \deleted{\unsure{donc le|donc ledit}} complexe de chaines $\mathbb{C}_{.}(X_*)_*$ peut être considéré comme ``le'' invariant additif universel'' de $X_*$. Quand on ``oublie $X_*$ au profit de $\mathbb{C}_{.}(X_*)^{\unsure{*}}=C_{.}\in \mathrm{Ab}_{.}$'', les foncteurs envisagés (en tant que foncteurs en $C_{.}$) sont encore exactement les foncteurs représentables sur $\mathrm{Ab}_{.}$.

[ \uline{NB}. Quand le foncteur $F$ a une structure plus riche, p.ex. une multiplication $F\otimes F\mathbin{\unsure{-|\to}}F$ i.e. est à valeurs dans les $\mathbb Z$-algèbres, \deleted{\unsure{est-ce|est $F($|et $F($}} cette structure sur $F(X_*)$ n'est pas reconstruite par
